\documentclass[10pt,a4paper]{amsart}
\usepackage[margin=19mm]{geometry}
\usepackage{amsmath,amssymb,mathtools}
\usepackage[scr=rsfs,cal=euler]{mathalfa}
\usepackage{xfrac}
\usepackage[unicode,hidelinks,bookmarks=false]{hyperref}
\newcommand{\ie}{i.e.\ }
\newcommand{\cf}{cf.\ }
\newcommand{\resp}{respectively\ }
\newcommand{\Subsection}{\subsection}
\providecommand{\nobreakdash}{\nobreak\hbox{-}}
\setlength{\emergencystretch}{2em}
\multlinegap0pt
\usepackage{amssymb}
\usepackage{float}
\usepackage{tikz}
\usepackage{enumerate}
\newtheorem{proposition}{Proposition}
\newtheorem{theorem}{Theorem}
\title{Arithmetic Hirzebruch-Zagier divisors and central derivative values of Rankin-Selberg $L$-functions}
\author{Jeanine Van Order}
\date{Source: arXiv:2510.10303v2 (CC BY 4.0)}
\begin{document}
\maketitle

\noindent\emph{Source and licence.} Arithmetic Hirzebruch-Zagier divisors and central derivative values of Rankin-Selberg $L$-functions. Version arXiv:2510.10303v2 (CC BY 4.0), distributed by arXiv under the Creative Commons Attribution 4.0 International licence, http://creativecommons.org/licenses/by/4.0/, as recorded in the arXiv licence metadata for this exact version and shown on the abstract page https://arxiv.org/abs/2510.10303v2. Full licence notice: \texttt{licenses/arxiv-2510.10303-CC-BY-4.0.txt}.\par\smallskip

\noindent\textit{Editorial scope.} Counted unit: the article's theorem RSIP (source0.tex lines 2884--2911). The article's complete original proof of it is delivered with it (source0.tex lines 2913--2971, one \texttt{proof} environment of 59 lines). No mathematics is written, repaired or re-derived: the statement and the proof are the article's own lines, in the article's own notation, and no retained line is substituted or re-flowed.  Retained uncounted as the source-local context the proof needs: lines 1797--1798; lines 1808--1810; lines 1863--1876; lines 2276--2281; lines 2391--2398; lines 2595--2596; lines 2598--2599; lines 2634--2645; lines 2767--2792.  Nothing was omitted from any retained span, so the package carries no omission record.  No retained line contains a citation, so the wrapper carries no bibliography and no bibliography entry.  No retained line contains a figure environment, an image-inclusion command or drawing code, so no image is delivered and the package declares no asset dependency.  Editorial formatting only, in addition: an AMS \texttt{amsart} preamble that restates the article's own declarations and macros used by the retained lines, namely the environments \texttt{proposition}, \texttt{theorem} on separate counters, together with the standard packages those lines need.  The retained environments print in this document as Proposition 1, Theorem 1 in order of appearance, whereas the article prints the same items with the numbering of its own full document.  The complete native source of the article as distributed in the arXiv e-print for this exact version is bundled unchanged as \texttt{sources/186880/source0.tex}.
\begin{align}\label{E-classical} E_L(\tau, s; l) := \sum\limits_{ \mu \in L{\vee}/L} E_L(g_{\tau}, s; \phi_{\infty}^l \otimes \lambda_f({\bf{1}}_{\mu}))
&= \sum\limits_{ \gamma \in \Gamma_{\infty} \backslash \Gamma } \left[ \Im(\tau)^{\frac{s + 1 - l}{2}} {\bf{1}}_0 \right] \Big\vert_{l, \omega_{L}} \gamma. \end{align}

\begin{align}\label{E-CM} E_{L_0}(\tau, s; l) &:= \sum\limits_{ \mu \in L_0^{\vee}/L_0 } 
E_{L_0}(g_{\tau}, s; \phi_{\infty}^l \otimes \lambda_f({\bf{1}}_{\mu})) {\bf{1}}_{\mu}, 
\quad g_{\tau} = n(u)m(v) \in \operatorname{SL}_2({\bf{R}}), \tau = u + iv \in \mathfrak{H}. \end{align}

\begin{proposition}\label{LFE-CM}

Suppose $(L_0, Q_0) = (\mathfrak{a}, -Q_{\mathfrak{a}})$ for $\mathfrak{a} \subset \mathcal{O}_k$ a nonzero integral 
ideal of an imaginary quadratic field $k = k(V_0)$ of discriminant $d_k$ and odd quadratic Dirichlet character $\eta_k(\cdot) = \left( \frac{d_k}{\cdot} \right)$,
with $Q_{\mathfrak{a}}(\cdot) = {\bf{N}}_{k/{\bf{Q}}}(\cdot)/{\bf{N}} \mathfrak{a}$ the corresponding positive definite norm form. 
Let $E_{L_0}(\tau; s; l)$ denote either of the corresponding Eisenstein series of weight $l=\lbrace 1, -1\rbrace$ defined in $(\ref{E-CM})$ and $(\ref{E-classical})$ above. Writing
\begin{align*} \Lambda(s, \eta_k) 
&= \vert d_k \vert^{\frac{s}{2}} \Gamma_{\bf{R}} (s+1) L(s, \eta_k), \quad \Gamma_{\bf{R}}(s) := \pi^{-\frac{s}{2}} \Gamma(s)\end{align*}
to denote the completed Dirichlet $L$-function $L(s, \eta_k)$ of the character $\eta_k$, the completed Eisenstein series
\begin{align*} E_{L_0}^{\star}(\tau, s; l) &:= \Lambda(s+1, \eta_k) E_{L_0}(\tau, s; l) = \Lambda(s+1, \eta_k) E_{L_0}(\tau, s; l)\end{align*}
satisfies the odd, symmetric functional equation 
\begin{align*} E_{L_0}^{\star}(\tau,s: l) &= -E_{L_0}^{\star}(\tau, -s; l). \end{align*}

\end{proposition}

\begin{theorem}[Bruinier-Yang]\label{BY4.7} We have that
\begin{align*} \Phi(f, Z(V_0)) &= - \frac{2}{\operatorname{vol}(K_0)} \left(
\operatorname{CT} \langle \langle f^+(\tau), \theta_{L_0^{\perp}}(\tau) \otimes \mathcal{E}_{L_0}(\tau) \rangle \rangle
+ L'(0, g \times \theta_{L_0^{\perp}}) \right) \\
&= - \frac{\operatorname{deg}(Z(V_0))}{2} \left( \operatorname{CT} \langle \langle f^+(\tau), \theta_{L_0^{\perp}}(\tau) \otimes \mathcal{E}_{L_0}(\tau) \rangle \rangle
+ L'(0, g \times \theta_{L_0^{\perp}}) \right). \end{align*} \end{theorem}

\begin{theorem}\label{realquad} 

Assume the signature $(n-1, 1)$ space $(W^{\perp}, Q) = (W^{\perp}, Q\vert_{W^{\perp}})$ is anisotropic. We have that
\begin{align*} \Phi(f, \mathcal{G}(W)) &= - \frac{2}{ \operatorname{vol}(K_W)} \left(
\operatorname{CT} \langle \langle f^+(\tau), {\bf{1}}_{L_W^{\perp} + 0} \otimes \mathcal{E}_{L_W}(\tau) \rangle \rangle
+ L'(0, g \times \theta_{L_W^{\perp}}) +I'(0, f \times \xi_{\frac{n-2}{2}}(\theta_{L_W^{\perp}})) \right). \end{align*} 

\end{theorem}

\begin{equation}\begin{aligned}\label{Hecketheta} \theta(\chi) \in \begin{cases} M_{1}(\Gamma_0(\vert d_k \vert), \eta_k) &\text{if $k$ imaginary quadratic}\\
A_0(\Gamma_0(\vert d_k \vert), \eta_k) &\text{if $k$ real quadratic}.\end{cases}\end{aligned}\end{equation}

\begin{align}\label{weight} l(k) := \begin{cases} 1 &\text{ if $k$ is imaginary quadratic} \\
0 &\text{ if $k$ is real quadratic} \end{cases}. \end{align}

\begin{proposition}\label{FE}

Let $\phi \in S_{l(\phi)}(\Gamma_0(N))$ be a normalized newform.
Assume $(N, d_k)=1$ and $l(\phi) \geq l(k)$, where $l(k) = \lbrace 0, 1 \rbrace$ as in $(\ref{weight})$ denotes the weight of the 
Hecke theta series $\theta(\chi)$ defined in $(\ref{Hecketheta})$. Put 
\begin{align*} L_{\infty}(s, \phi \times \theta(\chi)) &= \Gamma_{\bf{C}} \left( s +  \frac{l(\phi) -1}{2}  \right)^2
= \Gamma_{\bf{R}} \left( s +  \frac{l(\phi) -1}{2}  \right)^2 \Gamma_{\bf{R}}\left( s + \frac{l(\phi) +1}{2} \right)^2 \end{align*}
Then, the completed $L$-function 
\begin{align*} \Lambda(s, \phi \times \theta(\chi)) := L_{\infty}(s, \phi \times \theta(\chi)) L(s, \phi \times \theta(\chi)) \end{align*}
satisfies the symmetric functional equation 
\begin{align*} \Lambda(s, \phi \times \theta(\chi)) 
&= \eta_k(-N) \vert d_k N \vert^{1-2s}  \Lambda(1-s, \phi \times \theta(\chi)). \end{align*} \end{proposition}

\begin{proposition}\label{RS-equiv} 

Fix a holomorphic cuspidal newform $\phi \in S^{\operatorname{new}}_2(\Gamma_0(N))$ of weight $2$, level $\Gamma_0(N)$, and trivial character. 
Let $g_{\phi, A} \in S_2(\omega_{L_A}^{\vee})$ denote the lifting of $\phi$ to a vector-valued cusp form of weight $2$ 
and conjugate Weil representation $\omega_{L_A}^{\vee}$. We have the following identifications of completed Rankin-Selberg $L$-functions. \\

\begin{itemize}

\item[(i)] If $k$ is the imaginary quadratic field associated to the negative definite subspace $V_{A, 0} \subset V_A$
with $L_{A,0} = L_A \cap V_{A,0}$, we have the identifications of completed Rankin-Selberg $L$-functions 
\begin{align*} L^{\star}(2s-2, g_{\phi, A} \times \theta_{L_{A, 0}^{\perp}}) 
= \Lambda(s-1/2, \phi \times \theta_A) \end{align*} 
for each class $A \in C(\mathcal{O}_k)$, and for each class group character $\chi \in C(\mathcal{O}_k)^{\vee}$ the identification 
\begin{align*} \sum\limits_{A \in C(\mathcal{O}_k)} \chi(A) L^{\star}(2s-2, g_{\phi, A} \times \theta_{L_{A, 0}^{\perp}}) 
= \Lambda(s-1/2, \phi \times \theta(\chi)). \end{align*} 

\item[(ii)] If $k$ is the real quadratic field associated to the Lorentzian subspace 
$W_A \subset V_A$ with $L_{A,W} = W_A \cap L_A$, we have the identifications of completed Rankin-Selberg $L$-functions 
\begin{align*} L^{\star}(2s-2, g_{\phi, A} \times \theta_{L_{A, W}^{\perp}}) &= \Lambda(s-1/2, \phi \times \theta_A)\end{align*} 
for each class $A \in C(\mathcal{O}_k)$, and for each class group character $\chi \in C(\mathcal{O}_k)^{\vee}$ that 
\begin{align*} \sum\limits_{A \in C(\mathcal{O}_k)} \chi(A) L^{\star}(2s-2, g_{\phi, A} \times \theta_{L_{A, W}^{\perp}}) 
&= \Lambda(s-1/2, \phi \times \theta(\chi)).\end{align*}

\end{itemize}

\end{proposition}

\begin{theorem}\label{RSIP}

We retain the setup of Proposition \ref{RS-equiv}. For each class $A \in C(\mathcal{O}_k)$, 
let $f_{0, A} \in H_0(\omega_{L_A})$ be any harmonic weak Maass form 
whose image under the antilinear differential operator $\xi_0$ equals $g_{\phi, A}$, so
\begin{align*} \xi_0(f_{0,A})(\tau) &= g_{\phi, A}(\tau). \end{align*}
We have the following integral presentations of completed Rankin-Selberg $L$-functions, 
given in terms of sums over CM cycles or geodesic sets as in Theorems \ref{BY4.7} and \ref{realquad} above, respectively. \\

\begin{itemize}

\item[(i)] If $k$ is imaginary quadratic and $\eta_k(-N) = -\eta_k(N) = -1$, then 
\begin{align*} &\Lambda'(1/2, \phi \times \theta(\chi)) \\ &= - \Lambda(1, \eta_k) \sum\limits_{A \in C(\mathcal{O}_k)} \chi(A) 
\left[ \left( \frac{\operatorname{vol}(K_{A,0})  }{ 2} \right) \Phi(f_{0, A}, Z(V_{A, 0})) 
+ \operatorname{CT} \langle \langle f_{0, A}^+(\tau), \theta_{L_{A, 0}^{\perp}}(\tau) \otimes \mathcal{E}_{L_{A, 0}}(\tau) \rangle \rangle \right] \\
&= - \Lambda(1, \eta_k) \sum\limits_{A \in C(\mathcal{O}_k)} \chi(A) 
\left[ \frac{ \operatorname{deg}(Z(V_{A, 0}))}{2}  \Phi(f_{0, A}, Z(V_{A, 0})) 
+ \operatorname{CT} \langle \langle f_{0, A}^+(\tau), \theta_{L_{A, 0}^{\perp}}(\tau) \otimes \mathcal{E}_{L_{A, 0}}(\tau) \rangle \rangle \right]. \end{align*}

\item[(ii)] If $k$ is real quadratic and $\eta_k(-N) = \eta_k(N) = -1$, then 
\begin{align*} &\Lambda'(1/2, \phi \times \theta(\chi)) = - \Lambda(1, \eta_k) \\  &\times \sum\limits_{A \in C(\mathcal{O}_k)} \chi(A) 
\left[ \left( \frac{\operatorname{vol}(K_{A,W})  }{ 2} \right) \Phi(f_{0, A}, \mathcal{G}(W_A)) 
+ \operatorname{CT} \langle \langle f_{0, A}^+(\tau), \theta_{L_{A, W}^{\perp}}(\tau) \otimes \mathcal{E}_{L_{A, W}}(\tau) \rangle \rangle
+ I'(0, f_{0,A} \times \xi_0(\theta_{L_{A,W}^{\perp}})) \right]. \end{align*}

\end{itemize}

\end{theorem}

\begin{proof}

For (i), we have for each class $A \in C(\mathcal{O}_k)$ the relation
\begin{align*} \Phi(f_{0, A}, Z(V_{A, 0})) &= - \frac{2}{\operatorname{vol}(K_{A,0})} \left(
\operatorname{CT} \langle \langle f_{0, A}^+(\tau), \theta_{L_{A, 0}^{\perp}}(\tau) \otimes \mathcal{E}_{L_{A, 0}}(\tau) \rangle \rangle
+ L'(0, g_{\phi, A} \times \theta_{L_{A, 0}^{\perp}}) \right) \end{align*}
and hence 
\begin{align}\label{CMIP} - \left( \frac{ \operatorname{vol}(K_{A,0})  }{ 2} \right) \Phi(f_{0, A}, Z(V_{A, 0})) 
- \operatorname{CT} \langle \langle f_{0, A}^+(\tau), \theta_{L_{A, 0}^{\perp}}(\tau) \otimes \mathcal{E}_{L_{A, 0}}(\tau) \rangle \rangle 
&= L'(0, g_{\phi, A} \times \theta_{L_{A, 0}^{\perp}})  \end{align}
by Theorem \ref{BY4.7}. Observe that since 
$L^{\star}(s, g_{\phi, A} \times \theta_{L_{A,0}^{\perp}}) = - L^{\star}(s, g_{\phi, A} \times \theta_{L_{A,0}^{\perp}})$
by the odd, symmetric functional equation $E_{L_{A,0}}^{\star}(\tau, s; 1) = - E_{L_{A,0}}^{\star}(\tau, -s; 1)$ described in Proposition \ref{LFE-CM}, 
we have the vanishing of the central value 
$L^{\star}(0, g_{\phi, A} \times \theta_{L_{A,0}^{\perp}}) = L(0, g_{\phi, A} \times \theta_{L_{A,0}^{\perp}}) = 0$, and hence that 
\begin{align}\label{CCDV} \Lambda^{\star \prime}(0, g_{\phi, A} \times \theta_{L_{A,0}^{\perp}}) 
= \Lambda(1, \eta_k) L'(0, g_{\phi, A} \times \theta_{L_{A,0}^{\perp}}). \end{align} 
Moreover, observe that by the equivalence of $L$-functions shown in Proposition \ref{RS-equiv} (i)
with the functional equation for $\Lambda(s, \phi \times \theta(\chi))$ described in Proposition \ref{FE}, 
we are only in the non-degenerate situation when $\eta_k(-N) = -\eta_k(N) = -1$.
Hence, we can multiply each side of $(\ref{CMIP})$ to obtain the corresponding relation
\begin{equation} \begin{aligned}\label{CMIP2}  &- \Lambda(1, \eta_k) \left[ \left( \frac{ \operatorname{vol}(K_{A,0})  }{ 2} \right) \Phi(f_{0, A}, Z(V_{A, 0})) 
+ \operatorname{CT} \langle \langle f_{0, A}^+(\tau), \theta_{L_{A, 0}^{\perp}}(\tau) \otimes \mathcal{E}_{L_{A, 0}}(\tau) \rangle \rangle \right] 
= L^{\star \prime}(0, g_{\phi, A} \times \theta_{L_{A, 0}^{\perp}}). \end{aligned}\end{equation}
Taking a twisted linear combination of the $L$-values on each side of $(\ref{CMIP2})$ and using Proposition \ref{RS-equiv}, we obtain 
\begin{equation} \begin{aligned}\label{CMIP3} 
&- \Lambda(1, \eta_k) \sum\limits_{A \in C(\mathcal{O}_k)} \chi(A) \left[ \left( \frac{ \operatorname{vol}(K_{A,0})  }{ 2} \right) \Phi(f_{0, A}, Z(V_{A, 0})) 
+ \operatorname{CT} \langle \langle f_{0, A}^+(\tau), \theta_{L_{A, 0}^{\perp}}(\tau) \otimes \mathcal{E}_{L_{A, 0}}(\tau) \rangle \rangle \right] \\
&= \sum\limits_{A \in C(\mathcal{O}_k)} \chi(A) L^{\star \prime}(0, g_{\phi, A} \times \theta_{L_{A, 0}^{\perp}})
= \Lambda'(1/2, \phi \times \theta(\chi)) = \Lambda'(1/2, \phi \times \theta(\chi)). \end{aligned}\end{equation}

For (ii), we have for each class $A \in C(\mathcal{O}_k)$ the relation
\begin{align*} &\Phi(f_{0, A}, \mathcal{G}(W_A)) \\ &= - \frac{2}{ \operatorname{vol}(K_{A,W})} \left(
\operatorname{CT} \langle \langle f_{0, A}^+(\tau), \theta_{L_{A, W}^{\perp}}^+(\tau) \otimes \mathcal{E}_{L_{A, W}}(\tau) \rangle \rangle
+ L'(0, g_{\phi, A} \times \theta_{L_{A, W}^{\perp}}) +I'(0, f_{0,A} \times \xi_0(\theta_{L_{A,W}^{\perp}})) \right) \end{align*}
and hence 
\begin{equation}\begin{aligned}\label{geoIP} &- \left( \frac{ \operatorname{vol}(K_{A,W})  }{ 2} \right) \Phi(f_{0, A}, \mathcal{G}(W_A)) 
- \operatorname{CT} \langle \langle f_{0, A}^+(\tau), \theta^+_{L_{A, W}^{\perp}}(\tau) \otimes \mathcal{E}_{L_{A, W}}(\tau) \rangle \rangle 
- I'(0, f_{0,A} \times \xi_0(\theta_{L_{A,W}^{\perp}})) \\
&= L'(0, g_{\phi, A} \times \theta_{L_{A, W}^{\perp}})  \end{aligned}\end{equation}
by Theorem \ref{realquad}. Using the identification of completed $L$-functions of Proposition \ref{RS-equiv} (ii)
with the functional equation of Proposition \ref{FE}, we see that $L^{\star}(s, g_{\phi, A} \times \theta_{L_{A, W}^{\perp}})$
satisfies an odd symmetric functional equation when $\eta_k(-N) = \eta_k(N) = -1$, whence
$L^{\star}(0, g_{\phi, A} \times \theta_{L_{A,W}^{\perp}}) = L(0, g_{\phi, A} \times \theta_{L_{A,W}^{\perp}}) = 0$, and 
\begin{align*} \Lambda^{\star \prime}(0, g_{\phi, A} \times \theta_{L_{A,W}^{\perp}}) 
= \Lambda(1, \eta_k) L'(0, g_{\phi, A} \times \theta_{L_{A,W}^{\perp}}). \end{align*}
Hence, we can multiply each side of $(\ref{geoIP})$ to obtain the corresponding relation
\begin{equation} \begin{aligned}\label{geoIP2} 
&- \Lambda(1, \eta_k) \left[ \left( \frac{ \operatorname{vol}(K_{A,W})  }{ 2 } \right) \Phi(f_{0, A}, \mathcal{G}(W_A)) 
+ \operatorname{CT} \langle \langle f_{0, A}^+(\tau), \theta_{L_{A, W}^{\perp}}(\tau) \otimes \mathcal{E}_{L_{A, W}}(\tau) \rangle \rangle 
+ I'(0, f_{0,A} \times \xi_0(\theta_{L_{A,W}^{\perp}})) \right] \\
&= L^{\star \prime}(0, g_{\phi, A} \times \theta_{L_{A, W}^{\perp}}). \end{aligned}\end{equation}
Taking a twisted linear combination of the $L$-values on each side of $(\ref{geoIP2})$ and using Proposition \ref{RS-equiv}, we obtain 
\begin{equation} \begin{aligned}\label{geoIP3} 
&- \Lambda(1, \eta_k) \sum\limits_{A \in C(\mathcal{O}_k)} \chi(A) \left[ \left( \frac{ \operatorname{vol}(K_{A,W})  }{ 2 } \right) \Phi(f_{0, A}, \mathcal{G}(W_A)) 
+ \operatorname{CT} \langle \langle f_{0, A}^+(\tau), \theta_{L_{A, W}^{\perp}}(\tau) \otimes \mathcal{E}_{L_{A, W}}(\tau) \rangle \rangle 
+ I'(0, f_{0,A} \times \xi_0(\theta_{L_{A,W}^{\perp}})) \right] \\
&= \sum\limits_{A \in C(\mathcal{O}_k)} \chi(A) L^{\star \prime}(0, g_{\phi, A} \times \theta_{L_{A, W}^{\perp}})
= \Lambda'(1/2, \phi \times \theta(\chi)) = \Lambda'(1/2, \phi \times \theta(\chi)). \end{aligned}\end{equation} \end{proof}

\end{document}
